\section{Connecting NEAT and Graph Neural Networks}

NEAT and Graph Neural Networks emerged from different research traditions 




--- evolutionary computation and geometric deep learning respectively --- and have rarely been studied in conjunction. Nevertheless, a close examination of their computational structures reveals deep parallels that make their combination both natural and principled. This section reviews the literature that illuminates these connections and situates the present work within that context.

\subsection{Shared Computational Substrate: Graphs}

At the most fundamental level, both NEAT and GNNs operate on directed graphs. In NEAT, the genome directly encodes a directed graph: node genes represent neurons, and connection genes represent directed, weighted edges between them. The forward pass of a NEAT-evolved network is a graph computation --- information propagates from input nodes through hidden nodes to output nodes along the edges encoded in the genome, which need not form a simple layered structure and may include recurrent connections.

GNNs, by design, also perform computation over graphs: the message-passing procedure propagates information along the edges of the input graph, updating node representations iteratively. The architecture of a GNN --- the message function $M_k$, the update function $U_k$, and the number of layers $K$ --- determines how this graph computation unfolds.

Battaglia et al. \cite{battaglia2018relational} unified many neural architectures, including fully connected networks, convolutional networks, and recurrent networks, under the framework of \textit{Graph Networks}, demonstrating that these architectures can be seen as special cases of a general graph computation that operates on sets of nodes, edges, and a global context vector. This unification implies that the space of neural architectures over which NEAT searches already includes graph-network-like computations; evolving a NEAT genome is, in this view, a form of neural architecture search over the space of graph computations.

\subsection{Topology as Architecture}

In gradient-based deep learning, network topology is typically a fixed hyperparameter chosen prior to training. Architecture search methods such as Neural Architecture Search (NAS) \cite{zoph2017neural} attempt to automate this choice, but they generally operate over a restricted, predefined search space and use gradient surrogates or reinforcement learning rather than evolutionary exploration. NEAT, by contrast, treats topology as a first-class variable that evolves alongside the weights. Each genome in the NEAT population corresponds to a distinct graph topology, and the evolutionary process directly optimizes over this space.

For GNNs specifically, the topology of the message-passing graph is determined by the input data, not the architecture. However, the \textit{internal} computation graph of a GNN --- how information is transformed within each message and update function --- is itself a neural network that can vary in architecture. Evolving this internal network with NEAT is equivalent to performing open-ended architecture search over the GNN's computation graph, with the evolutionary objective provided by heuristic quality rather than gradient descent.

This perspective has been explored in the field of \textit{Evolved Graph Networks}: Guo et al. \cite{guo2021evolving} and related work have applied evolutionary search to discover GNN architectures for molecular property prediction, finding that evolved architectures outperform hand-designed ones on tasks with complex relational structure. Our work extends this line of research to the domain of combinatorial planning, where the fitness signal is the quality of the heuristic as measured by solve rate and search efficiency on Sokoban.

\subsection{Speciation and Modularity}

A key feature of NEAT is speciation, which protects structural innovations by grouping genomes according to topological similarity and applying fitness sharing within groups. This mechanism prevents premature convergence to simple topologies and allows novel structures time to refine their weights before being evaluated against incumbent networks. An analogous concept arises in GNN design through the notion of \textit{modularity}: complex GNN architectures are often built from modular components (encoders, message functions, aggregators, decoders) that can be independently designed and combined. NEAT's speciation can be seen as a population-level analogue of this modularity principle: distinct species correspond to distinct modular structures, and fitness sharing ensures that each structural hypothesis is evaluated on its own terms.

\subsection{Relational Inductive Biases}

Battaglia et al. \cite{battaglia2018relational} introduced the concept of \textit{relational inductive biases} to describe the structural assumptions encoded in different neural architectures about how entities and their relations determine outcomes. Convolutional networks encode a translational equivariance bias suited to grid-structured data; recurrent networks encode a sequential order bias suited to temporal data; GNNs encode a relational bias suited to graph-structured data.

NEAT, from this perspective, performs an evolutionary search that implicitly discovers the appropriate inductive biases for a given domain. By evolving networks that process graph-structured Sokoban states, NEAT can discover topologies that implement message-passing-like computations --- even without being explicitly constrained to do so --- if such computations provide a fitness advantage. This raises the intriguing hypothesis that NEAT applied to graph-structured inputs will converge to topologies that resemble GNN message-passing architectures, as these computations are naturally suited to relational reasoning over the board graph.

\subsection{Neuroevolution of GNNs for Planning}

The combination of neuroevolution and graph representations for planning problems has begun to attract attention in the literature. Alet et al. \cite{alet2019graph} proposed a modular meta-learning approach using graph neural networks for compositional reasoning, demonstrating that GNNs can generalize across combinatorial problem instances when trained on a distribution of graphs. The use of evolutionary methods to search over the space of such modular graph computations is a natural extension that our work explores.

For heuristic search specifically, Ferber et al. \cite{ferber2020neural} demonstrated that GNNs can learn admissible heuristics for planning problems in the STRIPS formalism by representing the planning domain as a graph. Their results show that the graph-structured representation enables generalization across problem instances of varying size and configuration --- a property that is difficult to achieve with flat input representations. Combining this insight with the topological search of NEAT provides a mechanism for discovering GNN architectures that are not only expressively suited to relational heuristic estimation but also structurally optimized for the specific distribution of Sokoban puzzles.

\subsection{Summary}

The parallel between NEAT and GNNs is not superficial. Both treat computation as a graph traversal, both exploit the relational structure of their respective inputs, and both can be understood as instances of the general graph computation framework of Battaglia et al. \cite{battaglia2018relational}. The primary difference is that NEAT evolves the computation graph (the network architecture) while GNNs operate on an externally provided input graph (the state representation). Combining them resolves this distinction: NEAT evolves the GNN's internal computation graph, and the GNN applies that evolved computation to the graph-structured Sokoban state. The result is a heuristic function that is both structurally expressive --- because it exploits relational structure --- and structurally adapted --- because its architecture was discovered by open-ended evolutionary search rather than fixed by human design.
